\section{Motivation}
\label{sec:motivation}

\paragraph{The problem.}
Preference-based post-training needs, for many prompts, a judgment of which of two responses is
better. Collecting these judgments from people is slow and expensive, so they are often produced by
another language model acting as a judge
\citep{bai2022constitutional,zheng2023judging,lee2024rlaif}. Judges are known to differ from human
raters in systematic ways, for example in their sensitivity to the order in which the two responses
are shown and to response length \citep{zheng2023judging}. A model trained on judge labels
therefore learns the judge's preferences, not those of people.

In practice some human labels are usually available, but only for a small fraction of the
comparisons. This leads to the question studied here: given judge outputs on all comparisons and
human labels on a few, how should we estimate what people would prefer on every comparison, and
how should a model be trained toward that estimate?

\paragraph{What the judge provides.}
A judge gives more than a verdict. We can read the probabilities it assigns to graded outcomes
(strong preference, weak preference, tie). We can also query it twice, with the responses shown in
both orders. This yields a win probability, and in addition a tie probability and a measure of how
much the judge's answer depends on the presentation order. These additional quantities may
indicate where the judge is unreliable relative to humans. A method that uses only the win
probability cannot exploit this.

\paragraph{Limitations of current approaches.}
Existing approaches fall into three groups.
\begin{itemize}
  \item \emph{Use the judge directly.} Judge probabilities serve as soft labels for reward modeling
        or preference optimization \citep{lee2024rlaif,furuta2024gapo}. This inherits the judge's
        systematic errors.
  \item \emph{Correct the judge with a few human labels.} This can be done with a fitted linear map
        \citep{vandenburg2025aligning}, a jointly weighted Bradley--Terry model \citep{cai2026dial},
        or by relabeling selected examples \citep{xu2025rlthf}. These methods either assume a fixed
        parametric relation between judge and human preferences or correct individual labels. They
        do not estimate how the judge's error varies across comparisons.
  \item \emph{Use the human labels alone.} This is statistically inefficient when labels are few.
\end{itemize}

\paragraph{Our approach.}
\citet{li2026personalizing} study a closely related problem in nonparametric regression. A
pretrained predictor is given, and a few labeled samples from the target population are used to
correct it. The difference between the target function and the predictor is estimated by a local
kernel estimator, after truncating the local variation of the predictor. The amount of truncation
is chosen by validation. The resulting estimator uses the predictor where it is accurate and
reduces to a standard kernel estimator where it is not, and it attains the minimax rate. Their
setting is scalar regression and does not involve preference data or model training.

We apply this method to LLM judges, with the judge's win probability in the role of the pretrained
predictor. The project has three parts, shown together in Figure~\ref{fig:overview}.
\begin{enumerate}
  \item \textbf{Estimation (\dfsp).} We estimate human preference probabilities by correcting the
        judge's win probability. The correction is estimated locally in the space of judge
        features (win probability, tie probability, order dependence).
  \item \textbf{Label selection.} We derive a rule for choosing which comparisons to send to human
        annotators. It minimizes a bound on the estimation error.
  \item \textbf{Post-training.} We use the estimated probabilities as targets for DPO-style
        training, and bound the trained model's preference risk by the estimation error.
\end{enumerate}

\begin{figure}[t]
\centering
\begin{tikzpicture}[
  node distance=5mm and 3.2mm,
  box/.style={draw, rounded corners=2pt, align=center, font=\footnotesize,
              minimum height=14mm, text width=23mm, inner sep=1.5pt},
  judge/.style={box, fill=MidnightBlue!8},
  human/.style={box, fill=BurntOrange!12},
  est/.style={box, fill=OliveGreen!12},
  arr/.style={-{Stealth[length=2mm]}, thick}
]
  % main row: judge side, estimation, post-training
  \node[box]                         (data)  {Comparisons\\ \mbox{$(x,y_A,y_B)$}\\ $N$ pairs};
  \node[judge, right=of data]        (judge) {LLM judge\\ 5 labels,\\ both orders};
  \node[judge, right=of judge]       (feat)  {Judge features\\ \mbox{$Z=(p_J,t_J,o_J)$}};
  \node[est,   right=of feat]        (dfsp)  {\dfsp\\ \mbox{$\hat p=p_J+\hat\delta$}\\ (Sec.~\ref{sec:dfsp})};
  \node[est,   right=of dfsp]        (dpo)   {Soft-label DPO\\ with targets $\hat p$\\ (Sec.~\ref{sec:post-training})};
  \node[box,   right=of dpo]         (model) {Post-trained\\ model};

  % human-label branch above
  \node[human, above=8mm of feat]    (sel)   {Label selection\\ (Sec.~\ref{sec:selection})};
  \node[human, above=8mm of dfsp]    (hum)   {Human labels\\ on \mbox{$n\ll N$} pairs};

  \draw[arr] (data)  -- (judge);
  \draw[arr] (judge) -- (feat);
  \draw[arr] (feat)  -- (dfsp);
  \draw[arr] (dfsp)  -- (dpo);
  \draw[arr] (dpo)   -- (model);
  \draw[arr] (feat)  -- (sel);
  \draw[arr] (sel)   -- (hum);
  \draw[arr] (hum)   -- (dfsp);
\end{tikzpicture}
\caption{Overview. The judge is run on all $N$ comparisons, and its output is summarized by the
features $Z$. Human labels are collected on $n$ comparisons chosen by the selection rule.
\dfsp{} corrects the judge's win probability using these labels, and the corrected probabilities
$\hat p$ are the training targets for the model.}
\label{fig:overview}
\end{figure}

